%% ****** Start of file apstemplate.tex ****** %
%%
%%
%%   This file is part of the APS files in the REVTeX 4.2 distribution.
%%   Version 4.2a of REVTeX, January, 2015
%%
%%
%%   Copyright (c) 2015 The American Physical Society.
%%
%%   See the REVTeX 4 README file for restrictions and more information.
%%
%
% This is a template for producing manuscripts for use with REVTEX 4.2
% Copy this file to another name and then work on that file.
% That way, you always have this original template file to use.
%
% Group addresses by affiliation; use superscriptaddress for long
% author lists, or if there are many overlapping affiliations.
% For Phys. Rev. appearance, change preprint to twocolumn.
% Choose pra, prb, prc, prd, pre, prl, prstab, prstper, or rmp for journal
%  Add 'draft' option to mark overfull boxes with black boxes
%  Add 'showkeys' option to make keywords appear
%\documentclass[aps,prl,preprint,groupedaddress]{revtex4-2}
\documentclass[aps,prl,twocolumn,groupedaddress]{revtex4-2}
%\documentclass[aps,prl,preprint,superscriptaddress]{revtex4-2}
%\documentclass[aps,prl,reprint,groupedaddress]{revtex4-2}

% You should use BibTeX and apsrev.bst for references
% Choosing a journal automatically selects the correct APS
% BibTeX style file (bst file), so only uncomment the line
% below if necessary.
%\bibliographystyle{apsrev4-2}

\usepackage{graphicx}
\usepackage{epstopdf}
\usepackage{caption}
%\usepackage{amsmath}% http://ctan.org/pkg/amsmath
%\usepackage{amsthm}
%\usepackage{amsfonts}
%\usepackage{subfigure}
%\usepackage{hhline}
%\usepackage[miktex]{gnuplottex}
%\usepackage{xcolor}
\usepackage{amssymb}
\usepackage{amsmath}
\usepackage{color}
\usepackage{hyperref}
%\usepackage[percent]{overpic}
\usepackage{tikz}
\usepackage{mathrsfs}
\usepackage{wasysym}
\usepackage{tikz-cd}
%\usepackage{linearb}
%\usepackage{wsuipa}
\usepackage{calc}
\usepackage{tipa}
\usepackage{booktabs}

%\usepackage{stix} %\fisheye
\usepackage{stackengine,scalerel}

% so sections, subsections, etc. become numerated.
\setcounter{secnumdepth}{3}

\newcommand{\FF}{\mathbb{F}}
\newcommand{\NN}{\mathbb{N}}
\newcommand{\ZZ}{\mathbb{Z}}
\newcommand{\QQ}{\mathbb{Q}}
\newcommand{\RR}{\mathbb{R}}
\newcommand{\CC}{\mathbb{C}}
\newcommand{\II}{\mathbb{I}}
\newcommand{\GG}{\mathbb{G}}

\DeclareMathOperator*{\argmax}{arg\,max}
\DeclareMathOperator*{\argmin}{arg\,min}
\newcommand{\avrg}[1]{\left\langle #1 \right\rangle}
\newcommand{\nelta}{\bar{\delta}}
\newcommand{\bra}[1]{\left\langle #1\right|}
\newcommand{\ket}[1]{\left| #1 \right\rangle}
\newcommand{\sbra}[1]{\langle #1|}
\newcommand{\sket}[1]{| #1 \rangle}
\newcommand{\bek}[3]{\left\langle #1 \right| #2 \left| #3 \right\rangle}
\newcommand{\sbek}[3]{\langle #1 | #2 | #3 \rangle}
\newcommand{\braket}[2]{\left\langle #1 \middle| #2 \right\rangle}
\newcommand{\ketbra}[2]{\left| #1 \middle\rangle \middle\langle #2  \right|}
\newcommand{\sbraket}[2]{\langle #1 | #2 \rangle}
\newcommand{\sketbra}[2]{| #1 \rangle  \langle #2 |}
\newcommand{\norm}[1]{\left\lVert#1\right\rVert}
\newcommand{\snorm}[1]{\lVert#1\rVert}
\newcommand{\bvec}[1]{\boldsymbol{\mathsf{#1}}}
\newcommand{\bcov}[1]{\boldsymbol{#1}}
\newcommand{\bdua}[1]{\boldsymbol{\check{#1}}}
%\newcommand{\bdov}[1]{\boldsymbol{\breve{#1}}}
\newcommand{\bdov}[1]{\breve{#1}}
%\newcommand{\bten}[1]{\boldsymbol{\mathfrak{#1}}}
\newcommand{\bten}[1]{\boldsymbol{\mathfrak{#1}}}
\newcommand{\forany}{\tilde{\forall}}
\newcommand{\qed}{$\overset{\circ}{.}\;$}
\newcommand{\bigo}{\mathcal{O}}
\newcommand{\spn}{\mathrm{spn}\,}
\newcommand{\krn}{\mathrm{krn}\,}
\newcommand{\biy}{\leftrightarrow}
\newcommand{\ugh}{\mathbin{\textlinb{\BPwheel}}}

\newcommand\bigeye{\ensurestackMath{\stackinset{c}{}{c}{-.3pt}%
  {\bullet}{\scriptstyle\bigcirc}}}
\newcommand\eye{\scalerel*{\bigeye}{x}}
%\newcommand*{\fisheye}{%
%    \mathbin{%
%        \ooalign{$\circledcirc$\cr\hidewidth$\bullet$\hidewidth}%
%    }%
%}

% double angle-bracket
% https://tex.stackexchange.com/questions/79657/how-to-get-double-angle-bracket-without-using-mnsymbol-package
\makeatletter
\newsavebox{\@brx}
\newcommand{\llangle}[1][]{\savebox{\@brx}{\(\m@th{#1\langle}\)}%
  \mathopen{\copy\@brx\mkern2mu\kern-0.9\wd\@brx\usebox{\@brx}}}
\newcommand{\rrangle}[1][]{\savebox{\@brx}{\(\m@th{#1\rangle}\)}%
  \mathclose{\copy\@brx\mkern2mu\kern-0.9\wd\@brx\usebox{\@brx}}}
\makeatother

% https://tex.stackexchange.com/questions/297362/bar-through-letter-in-mathmode
%\usepackage{calc}
\newsavebox\CBox
\newcommand\hcancel[2][0.5pt]{%
  \ifmmode\sbox\CBox{$#2$}\else\sbox\CBox{#2}\fi%
  \makebox[0pt][l]{\usebox\CBox}%  
  \rule[0.77\ht\CBox-#1/2]{\wd\CBox}{#1}}
\newcommand{\cb}{\hcancel{b}}
\newcommand{\cd}{\hcancel{d}}
\newcommand{\ct}{\hcancel{\partial}}
%\newcommand{\supp}{\mathsf{supp}}
\newcommand{\supp}{\mathrm{supp}}
\newcommand{\csupp}{\overline{\mathrm{supp}}}
\newcommand{\supr}[2]{\underset{#1}{\text{supr}}}

\begin{document}

% Use the \preprint command to place your local institutional report
% number in the upper righthand corner of the title page in preprint mode.
% Multiple \preprint commands are allowed.
% Use the 'preprintnumbers' class option to override journal defaults
% to display numbers if necessary
%\preprint{}

%Title of paper
\title{
M\'etodos Matem\'aticos de la F\'isica II\\
Cheat Sheet: Solving PDEs in Mathematical Physics
}

% repeat the \author .. \affiliation  etc. as needed
% \email, \thanks, \homepage, \altaffiliation all apply to the current
% author. Explanatory text should go in the []'s, actual e-mail
% address or url should go in the {}'s for \email and \homepage.
% Please use the appropriate macro foreach each type of information

% \affiliation command applies to all authors since the last
% \affiliation command. The \affiliation command should follow the
% other information
% \affiliation can be followed by \email, \homepage, \thanks as well.
\author{Juan I. Perotti}
\email[]{juan.perotti@unc.edu.ar}
%\homepage[]{Your web page}
%\thanks{}
%\altaffiliation{}
%\affiliation{}
\affiliation{Instituto de F\'isica Enrique Gaviola (IFEG-CONICET), Ciudad Universitaria, 5000 C\'ordoba, Argentina}
\affiliation{Facultad de Matem\'atica, Astronom\'ia, F\'isica y Computaci\'on, Universidad Nacional de C\'ordoba, Ciudad Universitaria, 5000 C\'ordoba, Argentina}

% \author{Colaborador}
% \affiliation{Instituto de F\'isica Enrique Gaviola (IFEG-CONICET), Ciudad Universitaria, 5000 C\'ordoba, Argentina}
% \affiliation{Facultad de Matem\'atica, Astronom\'ia, F\'isica y Computaci\'on, Universidad Nacional de C\'ordoba, Ciudad Universitaria, 5000 C\'ordoba, Argentina}
% \email[E-mail: ]{xxx.yyy@unc.edu.ar }

%Collaboration name if desired (requires use of superscriptaddress
%option in \documentclass). \noaffiliation is required (may also be
%used with the \author command).
%\collaboration can be followed by \email, \homepage, \thanks as well.
%\collaboration{Claudio J. Tessone}
%\noaffiliation

\date{\today}

%\begin{abstract}
%Bla bla bla... resumen...
%\end{abstract}

% insert suggested keywords - APS authors don't need to do this
%\keywords{}

%\maketitle must follow title, authors, abstract, and keywords
\maketitle
\onecolumngrid

% https://chatgpt.com/share/6a3bfb0e-e684-83e9-860b-2d90353eec33

\section{The General Philosophy}

Most linear PDEs encountered in physics can be written schematically as
$$
(\mathcal M u)(\mathbf x,t)
=
(\mathcal L u)(\mathbf x,t)
+
f(\mathbf x,t),
$$
where:
\begin{itemize}
\item $\mathcal L$ is a spatial differential operator,
\item $\mathcal M$ is a temporal differential operator,
\item $f$ is a source term.
\end{itemize}
The central idea is:
$$
\boxed{
\text{Geometry}
\rightarrow
\text{Boundary Conditions}
\rightarrow
\text{Spatial Eigenvalue Problem}
\rightarrow
\text{Eigenfunctions}
\rightarrow
\text{Expansion of Solutions}
}
$$
The fundamental step is therefore not solving the PDE directly, but solving the associated spatial eigenvalue problem.

\section{Spatial Eigenvalue Problem}

Given a spatial operator $\mathcal L$, solve the equation
$$
\mathcal L \phi_n
=
-\lambda_n \phi_n
$$
subject to the prescribed boundary conditions.
This reduces to the Helmholtz equation in the particular case where $\mathcal L=\nabla^2$.


The eigenfunctions ${\phi_n}$:
\begin{itemize}
\item form an orthogonal basis,
\item determine the structure of the solution, and
\item are reused for Laplace, Poisson, Heat, Wave, Helmholtz and Schr\"odinger equations.
\end{itemize}

\section{Sturm-Liouville Theory}

Most one-dimensional eigenvalue problems can be written as a Sturm-Liouville problem, which is about solving the ODE
$$
-\frac{d}{dx}
\left(
p\frac{dy}{dx}
\right)
+
q\,y
=
\lambda w\,y.
$$
under prescribed boundary conditions, and where $p:x\to p(x)\in \RR$, $q:x\to w(x)\in \RR$ and $w:x\to w(x)\in \RR$ denote some prescribed functions and $y:x\to y(x)\in \RR$ denotes a solution.

Properties:
\begin{itemize}
\item {\bf Real eigenvalues:}
$$
\lambda_n\in\mathbb R.
$$

\item {\bf Orthogonality:}
Two eigenfunctions $y_n$ and $y_m$ obtained for different eigenvalues $\lambda_n$ and $\lambda_m$ (i.e. for $n\neq m$) satisfy the orthogonality condition
$$
\langle y_n,y_m \rangle
=
0
$$
where
$$
\langle f,g \rangle
=
\int_a^b
f^*(x)g(x)w(x)\,dx
$$
denotes an inner product.

\item {\bf Completeness:}
Any sufficiently regular function can be expanded as
$$
f(x)
=
\sum_n
a_n y_n(x)
$$
where
$$
a_n 
= 
\frac{\langle y_n,f\rangle}{\langle y_n,y_n\rangle}
%=
%\int_a^b
%y_n^*(x)f(x)w(x)\,dx
$$
\end{itemize}

\section{Boundary Conditions}

Assume that $\Omega$ is the spatial domain under consideration.

\subsection{Dirichlet}

Value prescribed:
$$
u(\mathbf x,t)=g(\mathbf x,t)
\quad\text{for all $\mathbf x\in \partial\Omega$ and all $t$.}
$$
Examples:
\begin{itemize}
\item fixed temperature,
\item fixed electrostatic potential.
\end{itemize}

\subsection{Neumann}

Normal derivative prescribed:
$$
\frac{\partial u}{\partial \hat n}(\mathbf x,t)=g(\mathbf x,t)
\quad\text{for all $\mathbf x\in \partial\Omega$ and all $t$.}
$$
where $\hat n$ is the unit vector (pointing to the outside) normal to the boundary at each point $\mathbf x\in \partial \Omega$.
In other words, it is the unit vector orthogonal to the tanget space of $\partial \Omega$ at each point.

Examples:
\begin{itemize}
\item heat flux,
\item normal electric field.
\end{itemize}

\subsection{Robin (Mixed)}

$$
a\, u(\mathbf x,t)+b\,\frac{\partial u}{\partial \hat n}(\mathbf x,t)=g(\mathbf x,t)
\quad\text{for all $\mathbf x\in \partial\Omega$ and all $t$.}
$$

Examples:
\begin{itemize}
\item convective heat transfer.
\end{itemize}

\subsection{Periodic}

$$
u(\mathbf x+L,t)=u(\mathbf x,t).
$$
Examples:
\begin{itemize}
\item rings,
\item crystals,
\item Fourier analysis.
\end{itemize}

\section{Separation of Variables}

Assume
$$
u(\mathbf x,t)
=
X(\mathbf x)T(t).
$$
Then
$$
\frac{\mathcal M T}{T}
=
\frac{\mathcal L X}{X}
=
-\lambda.
$$
This produces:

\begin{itemize}
\item A {\bf spatial problem}
$$
\mathcal L X
=
-\lambda X.
$$
\item and a {\bf temporal problem} determined by the PDE under study.
\end{itemize}
Thus every separable PDE ultimately reduces to solving the a few spatial eigenvalue problems, since the heat, wave, Schr\"odinger, Poisson, etc., all involve the same spatial operator $\mathcal L$, and therefore the same eigenfunctions $\phi_n$.

\section{Coordinate Systems and Eigenfunctions}

\subsection{Cartesian Coordinates}

\subsubsection{Interval $[0,L]$}

Eigenproblem:
$$
X''+\lambda X=0.
$$

\begin{itemize}
\item Dirichlet
$$
X(0)=X(L)=0.
$$
Eigenfunctions:
$$
X_n(x)
=
\sin\frac{n\pi x}{L}.
$$
Eigenvalues:
$$
\lambda_n
=
\left(\frac{n\pi}{L}\right)^2.
$$
\item Neumann
$$
X'(0)=X'(L)=0.
$$
Eigenfunctions:
$$
X_n(x)
=
\cos\frac{n\pi x}{L}.
$$
Eigenvalues:
$$
\lambda_n
=
\left(\frac{n\pi}{L}\right)^2.
$$
\end{itemize}

\subsubsection{Periodic}

$$
X(0)=X(L),
\qquad
X'(0)=X'(L).
$$
Eigenfunctions:
$$
e^{i2\pi nx/L}.
$$
Eigenvalues:
$$
\lambda_n
=
\left(\frac{2\pi n}{L}\right)^2.
$$

\subsection{Rectangular Domains}

$$
0<x<L_x,
\qquad
0<y<L_y.
$$
Dirichlet eigenfunctions:
$$
\phi_{nm}(x,y)
=
\sin\frac{n\pi x}{L_x}
\sin\frac{m\pi y}{L_y}.
$$
Eigenvalues:
$$
\lambda_{nm}
=
\left(\frac{n\pi}{L_x}\right)^2
+
\left(\frac{m\pi}{L_y}\right)^2.
$$

\subsection{Polar Coordinates}

$$
(r,\theta).
$$
Laplacian:
$$
\nabla^2
=
\frac1r
\frac{\partial}{\partial r}
\left(
r\frac{\partial}{\partial r}
\right)
+
\frac1{r^2}
\frac{\partial^2}{\partial\theta^2}.
$$
\begin{itemize}
\item Angular Part
$$
\Theta''+m^2\Theta=0.
$$
Eigenfunctions:
$$
e^{im\theta}.
$$
Eigenvalues:
$$
m^2.
$$

\item Radial Part
$$
r^2R''
+rR'
+
(k^2r^2-m^2)R
=
0.
$$
Solutions:
$$
J_m(kr),
\qquad
Y_m(kr).
$$
(Bessel functions)
\end{itemize}

\subsection{Disk of Radius $a$}

Regular eigenfunctions:
$$
J_m(k_{mn}r)e^{im\theta}.
$$
with
$$
J_m(k_{mn}a)=0.
$$
Eigenvalues:
$$
\lambda_{mn}
=
k_{mn}^2
\left(
\frac{j_{mn}}{a}
\right)^2.
$$
where $j_{mn}$ denotes the $n$-th positive zero of $J_m$.
 
\subsection{Cylindrical Coordinates}

$$
(r,\phi,z).
$$
Eigenfunctions:
$$
J_m(k_r r)
e^{im\phi}
e^{ik_z z}.
$$
Eigenvalues:
$$
\lambda
=
k_r^2+k_z^2.
$$
where the dependence of $\lambda$ on $m$ goes through the value $k_r$ that appears in the radial equation.

\subsection{Spherical Coordinates}

$$
(r,\theta,\phi).
$$

\begin{itemize}
\item Angular Part
$$
L^2Y
=
\ell(\ell+1)Y.
$$
Eigenfunctions:
$$
Y_\ell^m(\theta,\phi).
$$
(Spherical harmonics)
Eigenvalues:
$$
\ell(\ell+1).
$$

\item Radial Part
$$
r^2R''
+
2rR'
+
\left(
k^2r^2-\ell(\ell+1)
\right)
R
=
0.
$$
The solutions are:
$$
j_\ell(kr)
\qquad
y_\ell(kr)
$$
where 
$$
j_\ell(x) = \sqrt{\frac{\pi}{2 x}}J_{\ell+\frac{1}{2}}(x)
\qquad
y_\ell(x) = \sqrt{\frac{\pi}{2 x}}Y_{\ell+\frac{1}{2}}(x)
$$
are the spherical Bessel functions.

\item Sphere of Radius $a$

Eigenfunctions:
$$
j_\ell(k_{\ell n}r)
Y_\ell^m(\theta,\phi).
$$
with
$$
j_\ell(k_{\ell n}a)=0.
$$
Eigenvalues:
$$
\lambda_{\ell n}
=
k_{\ell n}^2.
$$
\end{itemize}

\section{Common Special Functions}

%| Function                       | Differential Equation                     |
%| ------------------------------ | ----------------------------------------- |
%| Bessel (J_\nu,Y_\nu)           | (x^2y''+xy'+(x^2-\nu^2)y=0)               |
%| Legendre (P_\ell)              | ((1-x^2)y''-2xy'+\ell(\ell+1)y=0)         |
%| Associated Legendre (P_\ell^m) | Angular equation in spherical coordinates |
%| Hermite (H_n)                  | Harmonic oscillator                       |
%| Laguerre (L_n^\alpha)          | Hydrogen atom radial equation             |

\begin{table}[h]
\centering
\caption{Special Functions and Differential Equations}
\label{tab:special_functions}
\begin{tabular}{ll}
\toprule
\textbf{Function} & \textbf{Differential Equation} \\
\midrule
Bessel ($J_\nu, Y_\nu$) & $x^2y''+xy'+(x^2-\nu^2)y=0$ \\ \addlinespace
Legendre ($P_\ell$) & $(1-x^2)y''-2xy'+\ell(\ell+1)y=0$ \\ \addlinespace
Associated Legendre ($P_\ell^m$) & Angular equation in spherical coordinates \\ \addlinespace
Hermite ($H_n$) & Harmonic oscillator \\ \addlinespace
Laguerre ($L_n^\alpha$) & Hydrogen atom radial equation \\
\bottomrule
\end{tabular}
\end{table}

\section{Solving PDEs Using the Eigenbasis}

Once
$$
\mathcal L\phi_n=-\lambda_n\phi_n
$$
is known, expand
$$
u
=
\sum_n a_n(t)\phi_n.
$$
The PDE reduces to equations for the coefficients $a_n(t)$.

\section{Laplace Equation}

$$
\nabla^2u=0.
$$
No source term.
Solutions are harmonic functions.
Usually solved directly by separation of variables.

\section{Poisson Equation}

$$
\nabla^2u=f.
$$
Expand
$$
f
=
\sum_n f_n\phi_n.
$$
Then
$$
-\lambda_n a_n=f_n.
$$
Therefore
$$
a_n
=
-\frac{f_n}{\lambda_n}.
$$
Hence
$$
u
=
-\sum_n
\frac{f_n}{\lambda_n}
\phi_n.
$$

\section{Helmholtz Equation}

$$
\nabla^2u+k^2u=0.
$$
Equivalent to
$$
\mathcal L u=-k^2u.
$$
This is precisely the spatial eigenvalue problem.

\section{Heat Equation}

$$
u_t
=
D\nabla^2u.
$$
Expand
$$
u
=
\sum_n a_n(t)\phi_n.
$$
Then
$$
a_n'
=
-D\lambda_n a_n.
$$
Solution:
$$
a_n(t)
=
a_n(0)
e^{-D\lambda_n t}.
$$
Therefore
$$
u(\mathbf x,t)
=
\sum_n
a_n(0)
e^{-D\lambda_n t}
\phi_n(\mathbf x).
$$

\section{Wave Equation}

$$
u_{tt}
=
c^2\nabla^2u.
$$
Expansion gives
$$
a_n''
+
c^2\lambda_n a_n
=
0.
$$
Solution:
$$
a_n(t)
=
A_n\cos(c\sqrt{\lambda_n}t)
+
B_n\sin(c\sqrt{\lambda_n}t).
$$
Hence
$$
u
=
\sum_n
\left[
A_n\cos(c\sqrt{\lambda_n}t)
+
B_n\sin(c\sqrt{\lambda_n}t)
\right]
\phi_n.
$$

\section{Schrödinger Equation}

$$
i\hbar\frac{\partial\psi}{\partial t}
=
-\frac{\hbar^2}{2m}
\nabla^2\psi.
$$
Expansion yields
$$
a_n'
=
-i
\frac{\hbar\lambda_n}{2m}
a_n.
$$
Thus
$$
a_n(t)
=
a_n(0)
e^{-iE_nt/\hbar},
$$
where
$$
E_n
=
\frac{\hbar^2\lambda_n}{2m}.
$$

\section{Green Functions}

For
$$
\mathcal L u=f,
$$
define
$$
\mathcal L G(x,x')
=
\delta(x-x').
$$
Then
$$
u(x)
=
\int
G(x,x')
f(x')\,dx'.
$$

\subsection{Spectral Representation}

If
$$
\mathcal L\phi_n
=
-\lambda_n\phi_n,
$$
then
$$
G(x,x')
=
-\sum_n
\frac{\phi_n(x)\phi_n(x')}
{\lambda_n}.
$$
Thus Green functions are simply the inverse operator written in the eigenbasis.

{\bf Observation:}
Previous expansion assumes that all eigenvalues satisfy $\lambda_n\neq 0$.
If a zero mode exists (e.g. for Neumann or periodic boundary conditions), then the operator $\mathcal L$ is not invertible and the Green function must be modified by excluding the zero mode:
$$
G(x,x')
=
-\sum_{\lambda_n\neq0}
\frac{\phi_n(x)\phi_n(x')}
{\lambda_n}.
$$
In this case, solutions exist only if the source $f$ is orthogonal to the kernel of $\mathcal L$ (compatibility condition).

\section{Fourier Methods}

Useful on infinite domains.
Transform
$$
\hat u(k)
=
\int
u(x)e^{-ikx}\,dx.
$$
Derivatives become
$$
\partial_x
\rightarrow
ik,
$$
$$
\partial_x^2
\rightarrow
-k^2.
$$
Differential equations become algebraic equations in Fourier space.

\section{Master Formula}

Almost every linear PDE solved in a standard Mathematical Methods course eventually takes the form

$$
u(\mathbf x,t)
=
\sum_n
a_n(t)\phi_n(\mathbf x),
$$
where:
\begin{itemize}
\item $\phi_n$ are eigenfunctions of the spatial operator,
\item $\lambda_n$ are the corresponding eigenvalues,
\item $a_n(t)$ are determined by the temporal equation,
\item coefficients are fixed by initial and boundary conditions.
\end{itemize}
The entire theory can therefore be summarized as
$$
\boxed{
\text{PDE}
\longrightarrow
\text{Eigenvalue Problem}
\longrightarrow
\text{Eigenfunctions}
\longrightarrow
\text{Series Expansion}
\longrightarrow
\text{Solution}
}
$$
which unifies Fourier series, Bessel functions, spherical harmonics, Sturm-Liouville theory, Green functions, and separation of variables into a single spectral framework.


\end{document}
%
% ****** End of file apstemplate.tex ******


